\begin{abstract}
Retrieval-augmented generation systems increasingly detect contradictions among retrieved passages and flag them to
the generator. We ask what such a flag must say, and how accurate its detector must be, for it to help against
misinformation. On fixed passages from five benchmarks, with seven generators from three model families, we vary the
detector and the format of its flag, in hypotheses pre-registered in stages (H1--H9 before any confirmatory data,
H10--H31 in eight later addenda). A perfectly accurate neutral flag (``these passages conflict'') does not reliably
help. Its effect depends on the answer policy and the model family; pooled over \neuMlK\ settings, allowing
for shared questions, it is \neuMlMu\ points \neuMlCI. Naming the false side helps, but on average no more than
an answer key naming the wrong answer (pooled difference \PoolTwentyThree\ points, equivalent within $\pm$3 on average). In a
randomised test, a flag naming the false side raises accuracy by \hThirtyOneAMin\ to \hThirtyOneAMax\ points and one naming the correct side lowers it by
\ahThirtyOneBMax\ to \ahThirtyOneBMin. Randomised flags on true pairs therefore help only above a break-even direction
precision, which rises with the generator's accuracy from \beMin\ to \beMax. Real detectors are selective, so this
threshold alone can misjudge them. Combined with each cell's coverage and no-flag accuracy, the follow rates from this
test account retrospectively for the flag's effect in \predEightN\ crossed detector--generator cells
($r=\predEightR$, mean error \predEightMae\ points; exploratory). What a
detector contributes is its verdict: passed as an answer key, it matches the flag (pooled, equivalent within $\pm$3). On genuine
conflicts, direction removes valid answers, and when conflicts are rare a real detector's false flags on clean
questions can outweigh its gains. We release the harness, pre-registrations and every run.
\end{abstract}
